\documentclass[11pt,a4paper]{article}
\usepackage[T1]{fontenc}
\usepackage[utf8]{inputenc}
\usepackage{lmodern}
\usepackage[a4paper,margin=2.4cm]{geometry}
\usepackage{amsmath,amssymb,amsthm}
\usepackage{parskip}
\title{\textbf{Mathematical notes}\\[0.3em]\large Ancillary material for
``The Last Human Gate''}
\author{Jeremy Canale}
\date{September 2026}
\newtheorem{lemma}{Lemma}
\begin{document}
\maketitle

These notes derive the algebraic statements used in the paper. Results are referred to by name, as in the manuscript. Nothing here is estimated from data.

\section{Workload ratio and break-even thresholds}

With common visits $\lambda$ and useful hours $K$, the labor account gives
\[
\rho=\frac{N_A}{N_H}=\frac{\lambda\bar h_A+B_A}{\lambda h_0+B_H}
=\frac{\bar h_A/h_0+b_A}{1+b_H},\qquad b_s=\frac{B_s}{\lambda h_0}.
\]
Write $h_E=\kappa\eta h_0$, $h_R=\mu h_S^H=\mu\sigma h_0=mh_0$ and $h_W=rh_0$. Then
$\bar h_A/h_0=q\kappa\eta+(1-q)\sigma\mu+r$. The baseline partition states
$q\kappa+(1-q)\sigma=1$, so $(1-q)\sigma=1-\phi$ with $\phi=q\kappa$. Substituting gives both forms of
the workload ratio:
\[
\rho=\frac{q\kappa\eta+(1-q)\sigma\mu+r+b_A}{1+b_H}=\frac{\phi\eta+(1-\phi)\mu+r+b_A}{1+b_H}.
\]

\paragraph{Exception-share threshold.} In overall-mean units, $\rho<1$ is
$qk+(1-q)m+r+b_A<1+b_H$, that is $q(k-m)<1+b_H-m-r-b_A$. For $k>m$ this is the exception-share threshold $q<q^{*}$;
for $k<m$ the inequality reverses; for $k=m$ the case share drops out. Solving the same equality for
$k$ at a given $q>0$ gives $k^{*}=m+(1+b_H-m-r-b_A)/q$, the line drawn in the frontier figure, once per route.

\paragraph{Exception-effort margin.} In group-relative units, $\rho<1$ is
$\phi\eta<1+b_H-(1-\phi)\mu-r-b_A$. For $\phi>0$,
\[
\eta^{*}=\frac{1+b_H-(1-\phi)\mu-r-b_A}{\phi}=\mu+\frac{1+b_H-\mu-r-b_A}{\phi},
\qquad
\frac{\partial\eta^{*}}{\partial\phi}=-\frac{1+b_H-\mu-r-b_A}{\phi^{2}} .
\]
The margin therefore falls with $\phi$ exactly when $\mu+r+b_A<1+b_H$, that is, when a process made
only of standard cases would save labor. In the synthetic example $m$ is held common, so
$(1-\phi)\mu=(1-q)m$ and the code evaluates $\eta^{*}=[1+b_H-(1-q)m-r-b_A]/\phi$.

\section{Selection identity}

Let $T\ge 0$ be baseline handling time with $\mathbb{E}[T]=h_0>0$, $J\in\{0,1\}$ the residual
indicator, $e(T)=\Pr(J=1\mid T)$ and $q=\mathbb{E}[e(T)]>0$. Iterated expectations give
$\mathbb{E}[TJ]=\mathbb{E}[Te(T)]=\operatorname{Cov}(T,e(T))+h_0q$. Dividing by $q$ yields
$\mathbb{E}[T\mid J=1]-h_0=\operatorname{Cov}(T,e(T))/q$; dividing by $h_0$ yields
$\phi=q+\operatorname{Cov}(T,e(T))/h_0$. Since $\kappa=\mathbb{E}[T\mid J=1]/h_0$, also $\phi=q\kappa$.
For independent copies $T,T'$,
$2\operatorname{Cov}(T,e(T))=\mathbb{E}[(T-T')(e(T)-e(T'))]\ge 0$ when $e$ is nondecreasing, because
the two factors then share a sign. The identities restrict first moments and one cross-moment only:
truncating a distribution to a narrow band of difficult cases raises the conditional mean and can
lower the conditional variance.

\section{Strategic gate model}

\paragraph{Setup.} Let $\varepsilon$ have distribution function $F$ and density $f$. With
$s=\xi+u+g$, approval probability is $p(s)=\Pr(s+\varepsilon\ge\tau)=1-F(\tau-s)$, so
$p'(s)=f(\tau-s)\ge 0$ and $p''(s)=-f'(\tau-s)$. The project maximizes
$\Pi(u,g)=Vp(s)-c_u(u)-c_g(g)-D(g,a)$ over $u,g\ge 0$.

\paragraph{First- and second-order conditions.} At an interior optimum,
\[
Vp'(s)-c_u'(u)=0,\qquad Vp'(s)-c_g'(g)-D_g(g,a)=0 .
\]
Put $t=Vp''(s)$, $\omega_u=c_u''(u)>0$ and $\omega_g=c_g''(g)+D_{gg}(g,a)>0$. The Hessian is
\[
H=\begin{pmatrix} t-\omega_u & t\\ t & t-\omega_g\end{pmatrix},\qquad
\det H=\omega_u\omega_g-t(\omega_u+\omega_g)=\Delta .
\]
A negative-definite Hessian is characterized by $t-\omega_u<0$ and $\Delta>0$.
These sufficient second-order conditions for a strict local maximum hold automatically when $t\le 0$.

\paragraph{Assurance intensity.} Differentiating the two conditions with respect to $a$,
\[
H\begin{pmatrix}du\\ dg\end{pmatrix}=\begin{pmatrix}0\\ D_{ga}\,da\end{pmatrix}
\quad\Longrightarrow\quad
\frac{\partial u^{*}}{\partial a}=-\frac{tD_{ga}}{\Delta},\qquad
\frac{\partial g^{*}}{\partial a}=\frac{(t-\omega_u)D_{ga}}{\Delta},\qquad
\frac{\partial s^{*}}{\partial a}=-\frac{\omega_uD_{ga}}{\Delta}.
\]
With $D_{ga}>0$: cosmetic effort falls whenever the second-order conditions hold; the observable score
falls as well; genuine remediation rises weakly if and only if $t\le 0$. If $t>0$ (and the
second-order conditions hold), both efforts fall.

\begin{lemma}
If $f$ is differentiable and unimodal with mode zero, then $p''(s)\le 0$ whenever $s\ge\tau$.
The converse also holds if $f'(x)<0$ for every $x>0$.
\end{lemma}
\begin{proof}
Unimodality with mode zero means $f'(x)\ge 0$ for $x\le 0$ and $f'(x)\le 0$ for $x\ge 0$. Thus
$s\ge\tau$ gives $\tau-s\le 0$ and $p''(s)=-f'(\tau-s)\le 0$. Under the additional strict
derivative condition, $s<\tau$ gives $p''(s)>0$, proving the converse. Without that condition,
a flat density segment, including a zero-density tail, can give $p''(s)=0$ even below the threshold.
\end{proof}
The redirection result therefore applies to projects whose expected score already meets the
threshold; below it, stronger assurance can discourage both kinds of effort.

\paragraph{Corner solutions.} If $g^{*}=0$ because $Vp'(s)\le c_g'(0)+D_g(0,a)$, a marginal increase
in $a$ leaves $(u^{*},g^{*})$ unchanged. If $u^{*}=0$ is binding and $g^{*}>0$, the single remaining
condition gives $\partial g^{*}/\partial a=D_{ga}/(t-\omega_g)<0$ under its second-order condition
$t-\omega_g<0$.

\paragraph{Clearer criteria.} Let a clarity parameter $\theta$ lower both marginal costs:
$c'_{u\theta}<0$ and $c'_{g\theta}<0$. Then
\[
\frac{\partial u^{*}}{\partial\theta}=\frac{c'_{u\theta}(t-\omega_g)-t\,c'_{g\theta}}{\Delta},\qquad
\frac{\partial g^{*}}{\partial\theta}=\frac{(t-\omega_u)\,c'_{g\theta}-t\,c'_{u\theta}}{\Delta}.
\]
For $t\le 0$ the first term of each numerator is positive and the second is nonpositive, so neither
sign is determined: clearer criteria can raise genuine work, cosmetic work, or both.

\paragraph{Numerical check.} \texttt{test\_reproduce.py} solves the model with quadratic costs,
logistic noise and $D(g,a)=ag^{2}/2$ at two assurance levels and confirms the predicted signs in the
region $s>\tau$.

\section{Reliability bounds and the readiness example}

For events $\mathcal{C}_1,\dots,\mathcal{C}_n$ with $\Pr(\mathcal{C}_i)=p_i$, the Fr\'echet
inequalities give
$\max\bigl(0,1-\sum_i(1-p_i)\bigr)\le\Pr(\bigcap_i\mathcal{C}_i)\le\min_ip_i$. With $n=20$ and
$p_i=0.99$ the bounds are $0.80$ and $0.99$; independence gives $0.99^{20}\approx 0.8179$.

The geometric readiness index of $(0.01,1,1,1,1)$ is $0.01^{1/5}=10^{-0.4}\approx 0.398$, which shows
that the index is only partially sensitive to a near-zero dimension and motivates reporting a minimum
index and hard blockers alongside it.
\end{document}
